\providecolor{myRed}{RGB}{209,34,57}
\providecolor{myBlue}{RGB}{28,90,160}
\providecolor{myGreen}{RGB}{34,139,34}
\providecolor{myOrange}{RGB}{210,105,15}
\providecolor{statorIron}{RGB}{234,240,248}
\providecolor{rotorIron}{RGB}{246,247,250}
\providecolor{airGapCol}{RGB}{255,255,255}
\providecolor{fluxCol}{RGB}{20,80,165}
\providecolor{slotColor}{RGB}{248,242,235}

\begin{tikzpicture}[scale=1, >=Stealth, font=\small]

  % Geometry parameters
  \def\Rrotor{1.75}
  \def\Rgap{2.20}
  \def\Rstator{3.40}

  % 1. Iron Cores
  \fill[statorIron] (0,0) circle (\Rstator);
  \fill[airGapCol] (0,0) circle (\Rgap);
  \fill[rotorIron] (0,0) circle (\Rrotor);

  % Slot regions shading
  \fill[slotColor] (32:\Rgap) arc[start angle=32, end angle=148, radius=\Rgap] 
    -- (148:\Rgap+0.32) arc[start angle=148, end angle=32, radius=\Rgap+0.32] -- cycle;
  \fill[slotColor] (-32:\Rgap) arc[start angle=-32, end angle=-148, radius=\Rgap] 
    -- (-148:\Rgap+0.32) arc[start angle=-148, end angle=-32, radius=\Rgap+0.32] -- cycle;

  % Boundaries
  \draw[line width=1.1pt, black!75] (0,0) circle (\Rstator);
  \draw[line width=0.8pt, black!70] (0,0) circle (\Rgap);
  \draw[line width=0.8pt, black!70] (0,0) circle (\Rrotor);

  % Slot outlines
  \draw[line width=0.55pt, black!55] (32:\Rgap) -- (32:\Rgap+0.32) arc[start angle=32, end angle=148, radius=\Rgap+0.32] -- (148:\Rgap);
  \draw[line width=0.55pt, black!55] (-32:\Rgap) -- (-32:\Rgap+0.32) arc[start angle=-32, end angle=-148, radius=\Rgap+0.32] -- (-148:\Rgap);

  % 2. Magnetic Flux Lines
  % Center line along axis
  \draw[line width=0.6pt, fluxCol] (-\Rgap, 0) -- (\Rgap, 0);
  \draw[line width=0.6pt, fluxCol, ->] (0, 0) -- (0.08, 0);
  \draw[line width=0.6pt, fluxCol, ->] (\Rrotor, 0) -- ({(\Rrotor+\Rgap)/2 + 0.05}, 0);
  \draw[line width=0.6pt, fluxCol, ->] (-\Rgap, 0) -- ({-(\Rrotor+\Rgap)/2 + 0.05}, 0);

  % 7 nested symmetrical pairs covering the full depth of the stator yoke
  % dAng: air gap angle. Ri: yoke radius. phiArc: angle where circular arc starts
  \foreach \dAng/\Ri/\phiArc in {
    4/3.28/45,
    8/3.17/47,
    12/3.06/49,
    16/2.95/52,
    20/2.84/55,
    24/2.73/59,
    28/2.62/64
  } {
    \coordinate (RentryR) at (\dAng:\Rrotor);
    \coordinate (RgapR)   at (\dAng:\Rgap);
    \coordinate (RentryL) at ({180-\dAng}:\Rrotor);
    \coordinate (RgapL)   at ({180-\dAng}:\Rgap);

    % Upper half (y > 0)
    % A. Radial segment across air gap with arrow (strictly perpendicular to iron surface)
    \draw[line width=0.6pt, fluxCol, ->] (RentryR) -- (\dAng:{(\Rrotor+\Rgap)/2 + 0.04});
    \draw[line width=0.6pt, fluxCol] (\dAng:{(\Rrotor+\Rgap)/2}) -- (RgapR);

    \draw[line width=0.6pt, fluxCol, <-] (RentryL) -- ({180-\dAng}:{(\Rrotor+\Rgap)/2 - 0.04});
    \draw[line width=0.6pt, fluxCol] ({180-\dAng}:{(\Rrotor+\Rgap)/2}) -- (RgapL);

    % B. Inside rotor: smooth horizontal transition tangent to radial at boundary
    \pgfmathsetmacro{\yMid}{\Rrotor * sin(\dAng) * 0.95}
    \draw[line width=0.6pt, fluxCol]
      (RentryL) .. controls ($(RentryL) + ({-\dAng}:0.4)$) and (-0.5, \yMid) .. (0, \yMid)
      .. controls (0.5, \yMid) and ($(RentryR) + ({180+\dAng}:0.4)$) .. (RentryR);
    \draw[line width=0.6pt, fluxCol, ->] (-0.05, \yMid) -- (0.08, \yMid);

    % C. Stator yoke: fillet from radial at Rgap to azimuthal at Ri, true circular arc across top, and fillet back
    \coordinate (ArcStartU) at (\phiArc:\Ri);
    \coordinate (ArcEndU)   at ({180-\phiArc}:\Ri);
    \pgfmathsetmacro{\filletL}{(\Ri - \Rgap) * 0.6}
    \pgfmathsetmacro{\filletT}{\Ri * sin(\phiArc - \dAng) * 0.6}

    \draw[line width=0.6pt, fluxCol]
      (RgapR) .. controls ($(RgapR) + (\dAng:\filletL)$) and ($(ArcStartU) + ({\phiArc-90}:\filletT)$) .. (ArcStartU)
      arc[start angle=\phiArc, end angle={180-\phiArc}, radius=\Ri]
      .. controls ($(ArcEndU) + ({180-\phiArc+90}:\filletT)$) and ($(RgapL) + ({180-\dAng}:\filletL)$) .. (RgapL);

    % Lower half (y < 0)
    \coordinate (RentryRL) at (-\dAng:\Rrotor);
    \coordinate (RgapRL)   at (-\dAng:\Rgap);
    \coordinate (RentryLL) at ({-180+\dAng}:\Rrotor);
    \coordinate (RgapLL)   at ({-180+\dAng}:\Rgap);

    \draw[line width=0.6pt, fluxCol, ->] (RentryRL) -- (-\dAng:{(\Rrotor+\Rgap)/2 + 0.04});
    \draw[line width=0.6pt, fluxCol] (-\dAng:{(\Rrotor+\Rgap)/2}) -- (RgapRL);

    \draw[line width=0.6pt, fluxCol, <-] (RentryLL) -- ({-180+\dAng}:{(\Rrotor+\Rgap)/2 - 0.04});
    \draw[line width=0.6pt, fluxCol] ({-180+\dAng}:{(\Rrotor+\Rgap)/2}) -- (RgapLL);

    \draw[line width=0.6pt, fluxCol]
      (RentryLL) .. controls ($(RentryLL) + (\dAng:0.4)$) and (-0.5, -\yMid) .. (0, -\yMid)
      .. controls (0.5, -\yMid) and ($(RentryRL) + ({-180-\dAng}:0.4)$) .. (RentryRL);
    \draw[line width=0.6pt, fluxCol, ->] (-0.05, -\yMid) -- (0.08, -\yMid);

    \coordinate (ArcStartL) at (-\phiArc:\Ri);
    \coordinate (ArcEndL)   at ({-180+\phiArc}:\Ri);

    \draw[line width=0.6pt, fluxCol]
      (RgapRL) .. controls ($(RgapRL) + (-\dAng:\filletL)$) and ($(ArcStartL) + ({-\phiArc+90}:\filletT)$) .. (ArcStartL)
      arc[start angle=-\phiArc, end angle={-180+\phiArc}, radius=\Ri]
      .. controls ($(ArcEndL) + ({-180+\phiArc-90}:\filletT)$) and ($(RgapLL) + ({-180+\dAng}:\filletL)$) .. (RgapLL);
  }

  % 3. Conductors in stator slots
  \foreach \a in {38, 49, 60, 71, 82, 98, 109, 120, 131, 142} {
    \coordinate (C) at (\a:{\Rgap+0.16});
    \fill[white] (C) circle (2.6pt);
    \draw[line width=0.55pt, myRed!85!black] (C) circle (2.6pt);
    \fill[myRed!90!black] (C) circle (1.0pt);
  }
  \foreach \a in {-38, -49, -60, -71, -82, -98, -109, -120, -131, -142} {
    \coordinate (C) at (\a:{\Rgap+0.16});
    \fill[white] (C) circle (2.6pt);
    \draw[line width=0.55pt, myBlue!85!black] (C) circle (2.6pt);
    \draw[line width=0.55pt, myBlue!85!black] ($(C)+(-1.4pt,-1.4pt)$) -- ($(C)+(1.4pt,1.4pt)$);
    \draw[line width=0.55pt, myBlue!85!black] ($(C)+(-1.4pt,1.4pt)$) -- ($(C)+(1.4pt,-1.4pt)$);
  }

  % Winding schematic end-turn outline
  \draw[line width=0.45pt, black!50] (28:\Rgap+0.16) [bend right=25] to (45:\Rgap+0.36) arc[start angle=45, end angle=135, radius=\Rgap+0.36] [bend right=25] to (152:\Rgap+0.16);
  \draw[line width=0.45pt, black!50] (-28:\Rgap+0.16) [bend left=25] to (-45:\Rgap+0.36) arc[start angle=-45, end angle=-135, radius=\Rgap+0.36] [bend left=25] to (-152:\Rgap+0.16);

  % 4. Winding Axis
  \draw[dashed, line width=0.55pt, black!65] (-3.5, 0) -- (3.5, 0) 
    node[right=1pt, font=\scriptsize, align=left] {axis of stator\\winding};

  % Point P at angle 50 deg
  \draw[dashed, line width=0.5pt, black!60] (0,0) -- (50:3.65);
  \fill[black] (50:{\Rgap+0.08}) circle (1.8pt) node[left=1.5pt, font=\scriptsize] {P};

  % Angle phi arc
  \draw[->, line width=0.5pt, dashed, black!75] (0:3.15) arc[start angle=0, end angle=50, radius=3.15];
  \node[font=\footnotesize] at (25:3.45) {$\varphi$};

  % 5. Annotations
  \draw[dashed, line width=0.5pt, black!65, ->] (-3.2, -1.2) node[left=1pt, font=\scriptsize] {rotor} -- (-1.1, -0.5);
  \draw[dashed, line width=0.5pt, black!65, ->] (3.2, -1.7) node[right=1pt, font=\tiny, align=left] {air gap\\[-2pt](width exaggerated)\\[-2pt]to other dimensions} -- (2.0, -0.6);
  \draw[dashed, line width=0.5pt, black!65, ->] (3.2, -0.8) node[right=1pt, font=\scriptsize] {stator} -- (2.7, -0.7);

\end{tikzpicture}
